% Trechos orientativos removidos da seção de resultados de main.tex.
% Guardado para consulta futura; não é incluído no artigo principal.

% -------------------------------------------------------------------------
\subsection{Proposed Result Cases}
\label{subsec:proposed_result_cases}

For a journal-quality assessment, the numerical results should be organized
progressively, starting from a balanced operating condition and then moving
to the complete three-phase unbalanced formulation. This sequence makes it
possible to isolate the voltage-support capability of the compensators before
evaluating the additional benefit of phase-wise control. The following cases
are therefore recommended:

\begin{itemize}
    \item \textbf{Case A: Balanced system without compensation} --- reference
          condition used to identify the weak buses and the base loadability
          limit of the feeder;
    \item \textbf{Case B: Balanced system with fixed-angle SVCs} --- evaluates
          the benefit of conventional compensation without optimization;
    \item \textbf{Case C: Balanced system with optimized common firing angle}
          --- imposes $\alpha_{j,a}=\alpha_{j,b}=\alpha_{j,c}$ and separates
          the effect of optimized reactive-power support from phase balancing;
    \item \textbf{Case D: Unbalanced three-phase system without optimized
          control} --- identifies voltage-unbalance and harmonic-distortion
          problems under realistic phase loading;
    \item \textbf{Case E: Proposed phase-wise optimized control} --- uses
          independent values of $\alpha_{j,a}$, $\alpha_{j,b}$, and
          $\alpha_{j,c}$ at the controlled SVCs and represents the complete
          proposed methodology.
\end{itemize}

\noindent This organization is useful because Case C acts as an important
intermediate benchmark: if Case E outperforms Case C, the improvement can be
attributed to the phase-dependent control capability rather than only to the
presence of additional reactive-power compensation.

% =====================================================================

% -------------------------------------------------------------------------
\subsection{Three-Phase Unbalanced Loading Condition}
\label{subsec:unbalanced_results}

After the balanced study, the analysis should be extended to the full
three-phase unbalanced case. This is the most important part of the results
because it demonstrates the need for independent phase firing angles. The
unbalanced scenario should be created by applying different loading factors
to phases $a$, $b$, and $c$, especially at buses located downstream of the
controlled compensators. For example, phase $a$ may be stressed more heavily
than the other phases, while phase $c$ may be lightly loaded. The exact
percentages should be reported in a scenario table so that the study can be
reproduced.

\begin{table}[!t]
\centering
\caption{Suggested definition of three-phase loading scenarios}
\label{tab:unbalanced_scenarios_suggested}
\footnotesize
\renewcommand{\arraystretch}{1.2}
\begin{tabular}{@{}lccc@{}}
\toprule
\textbf{Scenario} & Phase $a$ & Phase $b$ & Phase $c$ \\
\midrule
Balanced & $1.00\lambda$ & $1.00\lambda$ & $1.00\lambda$ \\
Mild unbalance & -- & -- & -- \\
Severe unbalance & -- & -- & -- \\
Nonlinear-load case & -- & -- & -- \\
\bottomrule
\end{tabular}
\end{table}

\noindent The main result in this subsection should be a comparison between
common-angle control and phase-wise control. The quantities to emphasize are
the maximum VUF, the minimum phase voltage, and the phase-specific reactive
power injected by the SVCs. If the proposed method is effective, the
optimized angles should differ among phases and the most heavily loaded
phase should receive a different compensation level from the lightly loaded
phases.

\begin{table}[!t]
\centering
\caption{Suggested summary of three-phase unbalanced performance indices}
\label{tab:unbalanced_performance_suggested}
\footnotesize
\renewcommand{\arraystretch}{1.2}
\setlength{\tabcolsep}{2.5pt}
\begin{tabular}{@{}lcccc@{}}
\toprule
\textbf{Case} & $V_{\min}$ & VUF$_{\max}$ & THD$_{V,\max}$ & $\lambda_{\max}$ \\
 & (pu) & (\%) & (\%) & (pu) \\
\midrule
No optimized control & -- & -- & -- & -- \\
Common-angle opt. & -- & -- & -- & -- \\
Phase-wise opt. & -- & -- & -- & -- \\
\bottomrule
\end{tabular}
\end{table}

\begin{table}[!t]
\centering
\caption{Suggested reporting format for optimized firing angles}
\label{tab:optimized_angles_suggested}
\footnotesize
\renewcommand{\arraystretch}{1.2}
\begin{tabular}{@{}lcccc@{}}
\toprule
\textbf{SVC bus} & $\lambda$ & $\alpha_a$ & $\alpha_b$ & $\alpha_c$ \\
 &  & (deg) & (deg) & (deg) \\
\midrule
17 & -- & -- & -- & -- \\
32 & -- & -- & -- & -- \\
\bottomrule
\end{tabular}
\end{table}

\noindent Table~\ref{tab:optimized_angles_suggested} is essential for the
paper because it directly shows the action taken by the proposed controller.
For a strong discussion, the optimized angles should be plotted as functions
of $\lambda$ and compared with the voltage and VUF curves. This reveals
whether the optimizer changes the firing angles smoothly along the load curve
or whether abrupt transitions occur near active constraints.

% =====================================================================

% -------------------------------------------------------------------------
\subsection{Recommended Discussion Points}
\label{subsec:recommended_discussion_points}

The final part of the results section should synthesize the main findings in
terms of engineering interpretation. The discussion should emphasize four
points. First, in the balanced case, the optimized SVCs should improve the
voltage profile and increase the loading margin while maintaining nearly
equal phase firing angles. Second, in the unbalanced case, the proposed
phase-wise control should reduce VUF more effectively than common-angle
control because each phase can receive an independent reactive-power
contribution. Third, the harmonic results should confirm that THD constraints
remain satisfied at all operating points, avoiding a common drawback of
thyristor-based compensation. Fourth, the evolution of the firing angles
should show physically meaningful behavior: as the loading factor increases,
the controller should move the SVCs toward stronger capacitive support while
respecting the feasible range $90^\circ \leq \alpha_{j,\phi} \leq 180^\circ$.

For the paper to be competitive in a high-quality journal, the strongest
claim should be based on the simultaneous improvement of loadability,
voltage profile, VUF, and THD compliance. The results should therefore avoid
presenting only voltage profiles; instead, they should demonstrate that the
proposed phase-wise FC--TCR control provides a broader power-quality benefit
than fixed compensation or conventional balanced SVC control.